\chapter{Data Science e The Big Data Challenge}

\section{Introduzione ai Big Data}
Il concetto di \textbf{Big Data} ha molteplici definizioni, essendo un termine che si è evoluto rapidamente nel corso degli anni (spesso accompagnato da un forte \textit{hype} mediatico, come ironicamente illustrato in celebri vignette a tema). Tuttavia, da un punto di vista accademico e tecnologico, possiamo adottare la seguente definizione formale:

\begin{definition}[Big Data]
I Big Data sono dati la cui \textbf{scala, diversità e complessità} richiedono nuove architetture, tecniche, algoritmi e strumenti di analisi (\textit{analytics}) per poterli gestire ed estrarne valore e conoscenza nascosta.
\end{definition}

% FIGURA DA INSERIRE QUI
% Fonte: Slide 1
% Descrizione: Striscia a fumetti di Dilbert sull'hype dei Big Data.
\begin{figure}[H]
    \centering
    % \includegraphics[width=0.7\textwidth]{dilbert_big_data}
    \caption{L'hype mediatico attorno al concetto di Big Data.}
    \label{fig:dilbert_bigdata}
\end{figure}

\section{Le Sfide dei Big Data}
L'avvento dei Big Data introduce delle sfide significative che non possono essere affrontate con i sistemi tradizionali. Queste sfide si dividono in due macro-categorie:

\subsection{Tecnologia e infrastruttura}
La gestione di moli di dati così imponenti richiede nuovi paradigmi di programmazione e nuove architetture hardware/software. Il principio fondamentale che guida queste nuove tecnologie è:
\begin{center}
    \textit{"Transfer the processing power to the data"} \\
    (Trasferire la potenza di calcolo verso i dati)
\end{center}
Invece di spostare enormi quantità di dati verso i server di calcolo (operazione lenta e costosa in termini di rete), si spostano gli algoritmi e il calcolo direttamente sui nodi dove i dati sono immagazzinati. L'ecosistema di riferimento per questo paradigma è \textbf{Apache Hadoop} e \textbf{Apache Spark}.

% FIGURA DA INSERIRE QUI
% Fonte: Slide 4
% Descrizione: Logo dell'elefantino di Apache Hadoop.
\begin{figure}[H]
    \centering
    % \includegraphics[width=0.3\textwidth]{hadoop_elephant}
    \caption{L'ecosistema Apache Hadoop è una delle tecnologie fondanti per gestire le infrastrutture Big Data.}
    \label{fig:hadoop_elephant}
\end{figure}

\subsection{Gestione e analisi dei dati}
Oltre all'infrastruttura, c'è una nuova e forte enfasi sul "dato" in sé. Raccogliere i dati non basta: è necessario analizzarli per estrarne valore. Questa esigenza ha dato vita a una disciplina specifica: la \textbf{Data Science}.

\section{Data Science}
La Data Science è la disciplina che si occupa di \textit{"estrarre significato da grandissime quantità di dati"}. 
È un campo fortemente multidisciplinare che unisce statistica, machine learning, data mining, visualizzazione dei dati e gestione dei database.

\begin{remark}
Il termine \textit{"Data Scientist"} è stato coniato originariamente da D.J. Patil (ex Chief Data Scientist degli Stati Uniti), contribuendo a definire una nuova figura professionale fondamentale nel panorama tecnologico moderno.
\end{remark}

% FIGURA DA INSERIRE QUI
% Fonte: Slide 5
% Descrizione: Mappa concettuale che mostra le varie discipline che compongono la Data Science (Statistic, Machine Learning, Data Mining, ecc.).
\begin{figure}[H]
    \centering
    % \includegraphics[width=0.6\textwidth]{data_science_mindmap}
    \caption{La natura multidisciplinare della Data Science.}
    \label{fig:ds_mindmap}
\end{figure}

\section{Il Processo di Data Science (KDD)}
Il processo di estrazione della conoscenza dai dati è formalmente noto come \textbf{KDD} (\textit{Knowledge Discovery in Databases}). 

Il flusso di lavoro della Data Science si articola in una macro-pipeline composta da quattro fasi principali, disposte in sequenza logica:
\begin{enumerate}
    \item \textbf{Generation} (Generazione del dato)
    \item \textbf{Acquisition} (Acquisizione)
    \item \textbf{Storage} (Memorizzazione)
    \item \textbf{Analysis} (Analisi)
\end{enumerate}

% FIGURA DA INSERIRE QUI
% Fonte: Slide 6
% Descrizione: I due schemi del processo KDD. Da un lato il processo lineare (Data -> Preprocessing -> Data Mining -> Knowledge), dall'altro il ciclo di feedback continuo con il mondo reale (Real World -> Raw Data -> Model -> Data Product -> Real World).
\begin{figure}[H]
    \centering
    % \includegraphics[width=0.9\textwidth]{kdd_process}
    \caption{Rappresentazione del processo KDD (Knowledge Discovery in Databases) e del ciclo di vita di un Data Product.}
    \label{fig:kdd_process}
\end{figure}

\subsection{Fase 1: Generazione dei dati (Generation)}
La prima fase del processo KDD riguarda il modo in cui i dati vengono creati nel mondo reale. Possiamo classificare la generazione dei dati in tre tipologie principali:

\begin{itemize}
    \item \textbf{Registrazione passiva (Passive recording):} Produce tipicamente \textbf{dati strutturati}. Esempi classici includono le transazioni bancarie e di trading, i registri degli acquisti nei negozi (shopping records) e gli archivi della pubblica amministrazione (government sector archives). L'utente non genera attivamente il dato, ma compie un'azione che viene passivamente registrata dal sistema.
    
    \item \textbf{Generazione attiva (Active generation):} Produce tipicamente \textbf{dati semi-strutturati o non strutturati}. Si tratta principalmente di contenuti generati direttamente dagli utenti (\textit{User-generated content}), come ad esempio i post, le foto e i video caricati sui social network.
    
    \item \textbf{Produzione automatica (Automatic production):} Riguarda dati generati in modo continuo e automatico da dispositivi. Questi dati sono spesso \textit{location-aware} (geolocalizzati), \textit{context-dependent} (dipendenti dal contesto) e generati da dispositivi altamente mobili (come smartphone, sensori IoT, smartwatch).
\end{itemize}